\documentclass[11pt,a4paper]{article}
\usepackage[T1]{fontenc}
\usepackage[utf8]{inputenc}
\usepackage{mathptmx}
\AtBeginDocument{\let\hbar\relax
  \DeclareMathSymbol{\hbar}{\mathord}{AMSb}{"7E}}
\usepackage{amsmath,amssymb,amsthm}
\usepackage{geometry}
\usepackage{microtype}
\usepackage{hyperref}
\usepackage{booktabs}
\usepackage{enumitem}
\usepackage[numbers]{natbib}
\usepackage{titlesec}

\geometry{top=1.1in,bottom=1.1in,left=1.15in,right=1.15in}

\titleformat{\section}{\normalfont\bfseries\large}{\thesection}{1em}{}
\titleformat{\subsection}{\normalfont\bfseries\normalsize}{\thesubsection}{1em}{}
\titleformat{\subsubsection}{\normalfont\bfseries\itshape\normalsize}
  {\thesubsubsection}{1em}{}
\titlespacing*{\section}{0pt}{14pt}{4pt}
\titlespacing*{\subsection}{0pt}{10pt}{3pt}

\theoremstyle{plain}
\newtheorem{theorem}{Theorem}[section]
\newtheorem{proposition}[theorem]{Proposition}
\theoremstyle{definition}
\newtheorem{definition}[theorem]{Definition}
\newtheorem{prediction}[theorem]{Prediction}
\theoremstyle{remark}
\newtheorem{remark}[theorem]{Remark}

\hypersetup{colorlinks=true,linkcolor=black,citecolor=black,urlcolor=black,
  pdftitle={The Origin, Structure, and Fate of Universes},
  pdfauthor={Joshua F. Sandeman}}

\newcommand{\Pact}{\mathcal{P}_{\mathrm{act}}}

% ─────────────────────────────────────────────────────────────────────────────
\begin{document}

\begin{center}
  {\LARGE\bfseries The Origin, Structure, and Fate\\
   of Universes in Causal Graph Cosmology\par}
  \vspace{6pt}
  {\normalsize Paper~III of the Relational Actualism Series\par}
  \vspace{14pt}
  {\large Joshua F.\ Sandeman$^{1*}$\par}
  \vspace{4pt}
  {\normalsize $^{1}$Independent Researcher, Salem, OR, USA.\par}
  {\normalsize $^{*}$Corresponding author: sansuikyo@gmail.com
   $\cdot$ relationalactualism.org\par}
  \vspace{14pt}
\end{center}

\begin{abstract}
If physical reality is an irreversibly growing directed acyclic graph governed by
the unique BDG growth rule (Paper~I), what does this imply for cosmology?  We show
that the framework resolves five outstanding cosmological puzzles simultaneously:
(i)~the cosmological constant is identically zero, because unactualized processes
do not source the gravitational field; (ii)~dark matter is not a particle species
but the actualization-density topology of sparse regions where effective
dimensionality reduces from four to two, producing flat rotation curves from
baryonic matter alone; (iii)~the Hubble tension is a density-dependent expansion
effect, with a specific prediction $H_0 \approx 75.9$~km/s/Mpc in the Eridanus
supervoid; (iv)~the CMB anomalies---low quadrupole, axis-of-evil alignment,
hemisphere asymmetry---are signatures of the spinning supermassive black hole from
which our universe nucleated, diluted by a parameter-free $(2/\ell)^2$ law; and
(v)~the universe does not reach thermal equilibrium because void--filament dynamics
fragment every connected component before heat death, prohibiting Boltzmann Brains
dynamically.  We derive the parent black hole mass
($\approx 2.5$--$6.6 \times 10^6\,M_\odot$), the three-parameter CMB anomaly
model that accounts for five independent observations with zero remaining freedom,
and the generational cycle by which universes reproduce through causal severance.
\end{abstract}

%================================================================
\section{Introduction}\label{sec:intro}
%================================================================

Modern cosmology harbours deep puzzles at both ends of the energy scale.  The
cosmological constant~$\Lambda$ accounts for 68\% of the energy density, yet QFT
predicts a vacuum energy $10^{120}$ times larger.  Dark matter accounts for 27\%,
yet decades of direct detection have found nothing.  The Hubble constant measured
locally ($H_0 \approx 73$~km/s/Mpc) disagrees with the early-universe value
($H_0 \approx 67.4$~km/s/Mpc) at 4--5$\sigma$.  The largest CMB angular scales
display anomalies that $\Lambda$CDM cannot explain.  And the long-term fate, in the
standard picture, is thermal death.

This paper shows that all five puzzles are resolved by a single hypothesis: physical
reality is an irreversibly growing directed acyclic graph whose growth rule is
uniquely determined by self-consistency (Paper~I~\cite{P1}).  Each resolution is a
structural consequence, not a post-hoc accommodation.

%================================================================
\section{General Relativity and the Cosmological Constant}\label{sec:gr}
%================================================================

\subsection{GR from BDG uniqueness}\label{sec:gr-deriv}

The Einstein field equations are derived, not imposed.  The derivation uses three
Lean-verified results and one published theorem:

\begin{enumerate}[nosep]
\item The Local Ledger Condition (L01,~[LV]) gives local conservation at every
  vertex.
\item Amplitude locality (O01,~[LV]) constrains the growth functional to depend
  only on the causal past.
\item The BDG closure theorem (L11,~[LV]) gives $d = 4$.
\item The Benincasa--Dowker uniqueness theorem~\cite{BenincasaDowker2010}
  establishes the BDG action as the unique local causal-set action in $d = 4$,
  whose continuum limit is the Einstein--Hilbert action.
\end{enumerate}

Variation of the Einstein--Hilbert action gives the field equations:
\begin{equation}
G_{\mu\nu} = 8\pi G\,\Pact[T_{\mu\nu}],
\label{eq:efe}
\end{equation}
where $\Pact$ is the actualization projector.  The Bianchi identity
$\nabla_\mu G^{\mu\nu} = 0$ is not a separate geometric fact---it is the LLC
viewed at the level of the metric.  In standard GR, the geometric Bianchi identity
and matter conservation $\nabla_\mu T^{\mu\nu} = 0$ are mysteriously compatible.
In RA, they are the \emph{same} conservation law (the LLC) at different levels of
description.

\subsection{Structural zero cosmological constant}\label{sec:lambda}

The actualization projector removes all off-shell contributions from the
gravitational source.  Vacuum fluctuations have $\Delta S = 0$ and do not cross the
actualization threshold.  Therefore:
\begin{equation}
\Lambda = 0 \quad\text{(structural, not fine-tuned).}
\label{eq:lambda}
\end{equation}
The $10^{120}$ vacuum energy discrepancy dissolves: the QFT calculation sums over
virtual processes that never become vertices in the causal graph.

%================================================================
\section{Dark Matter as Actualization-Density Topology}\label{sec:dm}
%================================================================

\subsection{The WIMP prohibition}\label{sec:wimp}

Any particle that does not participate in on-shell electromagnetic or strong
interactions has zero RA assembly depth and contributes nothing to the gravitational
source.  WIMPs, axions, and sterile neutrinos are \emph{categorically prohibited}
as dark matter candidates.  Current data from LZ and XENONnT are consistent: no
signal above the neutrino floor.

\subsection{Dimensionality reduction in sparse regions}\label{sec:dim-red}

In regions of low actualization density, the causal graph becomes effectively
two-dimensional, following the Durhuus--Jonsson--Wheater
theorem~\cite{DJW2010} for Lorentzian random graphs.  In two-dimensional effective
geometry, the gravitational dynamics include a logarithmic correction that flattens
rotation curves without additional matter.  The boundary correction term is:
\begin{equation}
\nabla^2(\ln\lambda) = \text{local actualization density gradient},
\label{eq:rotation}
\end{equation}
where $\lambda$ is the local actualization density.  This reproduces observed flat
galactic rotation curves from baryonic matter alone.

\subsection{The baryon-to-dark-matter ratio}\label{sec:f0}

The ratio $\Omega_{\mathrm{CDM}}/\Omega_b$ is derived from BDG path weights at
$\mu = 1$ (Paper~I~\cite{P1}): $f_0 = 5.42$, in 0.3\% agreement with
Planck~2018 ($5.416 \pm 0.015$)~\cite{Planck2020}.  In RA, ``dark matter'' is the
actualization-density topology of non-baryonic regions---not a new particle, but the
same graph physics producing different effective geometry at different densities.

%================================================================
\section{The Hubble Tension}\label{sec:hubble}
%================================================================

The RA field equations couple the local expansion rate to baryon density: regions
with lower actualization density expand faster.  The predicted Hubble constant
within the KBC void is:
\begin{equation}
H(\text{Eridanus}) \approx 75.9~\text{km/s/Mpc} \quad[\text{CV}],
\label{eq:hubble}
\end{equation}
consistent with the locally measured high value.  The globally averaged value
remains near 67.4~km/s/Mpc.  The tension is not a discrepancy between
measurements---it is a real physical effect: different regions genuinely have
different expansion rates because they have different actualization densities.

\begin{prediction}[Hubble gradient]\label{pred:hubble}
$H_0$ correlates linearly with baryon density along the line of sight.
Testable with DESI.  The correlation should be absent in $\Lambda$CDM.
\end{prediction}

%================================================================
\section{Causal Severance and Black Holes}\label{sec:severance}
%================================================================

\subsection{Bandwidth saturation}\label{sec:bandwidth}

The causal graph has a maximum actualization rate---the Planck rate.  When a
collapsing region approaches this limit, $\mu$ climbs toward $\mu_2 \approx 5.7$,
where $P_{\mathrm{acc}} \to 0$.  The graph cannot grow further locally.

\subsection{Topological disconnection}\label{sec:disconnection}

When $P_{\mathrm{acc}} \to 0$, no new vertices are accepted.  The region severs:
no outgoing causal edges exist.  The graph cut theorem~[LV] guarantees the LLC
holds independently on each side.  Event horizons are boundaries at which the graph
has no outgoing edges.  Conservation is not violated across them; it is
\emph{undefined}, because there is no causal connection across which to define it.

\subsection{No singularities}\label{sec:no-sing}

The GR singularity is replaced by a causal termination: a region where $\mu$
exceeds $\mu_2$ and the graph stops growing.  There is no point of infinite density
because the graph is discrete---dividing by zero is a continuum artefact.

%================================================================
\section{The Parent Black Hole}\label{sec:parent}
%================================================================

At causal severance, the daughter inherits its baryon-to-photon ratio~$\eta_b$ from
the Bekenstein--Hawking entropy of the parent.  Our universe has
$\eta_b = 6.1 \times 10^{-10}$, requiring $S_{\mathrm{BH}}/k_B \approx 10^{90}$.
Inverting the Bekenstein--Hawking formula:
\begin{equation}
M_{\mathrm{parent}} \approx 2.5\text{--}6.6 \times 10^6\,M_\odot,
\label{eq:parent}
\end{equation}
a supermassive black hole comparable to Sagittarius~A$^*$.  Stellar-mass black
holes give the wrong~$\eta_b$; the SMBH mass range is uniquely determined.

%================================================================
\section{Nucleation: The Kerr Geometry}\label{sec:kerr}
%================================================================

\subsection{The parent's spin}\label{sec:spin}

A real SMBH spins, with $a_* \in [0,1)$.  Observed SMBHs in the relevant mass
range typically have $a_* \approx 0.7$--$0.95$.

\subsection{Angular decomposition of the boundary flux}\label{sec:angular}

The Kerr horizon metric function $\Sigma_+({\theta}) = r_+^2 + a^2\cos^2\theta$
decomposes into Legendre polynomials:
\begin{equation}
\Phi_{\mathrm{boundary}}(\theta) = \Phi_0\bigl[1 + \varepsilon_2\,P_2(\cos\theta)
+ \varepsilon_3\,P_3(\cos\theta)\bigr],
\label{eq:boundary}
\end{equation}
where $\varepsilon_2(a_*) \approx 0.036\,a_*^2 + 0.326\,a_*^4$ from the Kerr
geometry~[CV], and $\varepsilon_3 = 0$ \emph{identically} by equatorial symmetry.
Any octupole requires accretion asymmetry.  Both $\varepsilon_2$ and
$\varepsilon_3$ share the spin axis---the quadrupole from geometry, the octupole
from the accretion disk in the equatorial plane.

\subsection{Daughter initial conditions}\label{sec:daughter}

The daughter is born with anisotropic initial conditions:
\begin{equation}
\frac{\delta\rho}{\rho}(\theta,\phi) = \varepsilon_2\,P_2(\cos\theta)
+ \varepsilon_3\,P_3(\cos\theta),
\label{eq:ic}
\end{equation}
with the preferred axis equal to the parent's spin axis.

%================================================================
\section{The $(2/\ell)^2$ Dilution Law}\label{sec:dilution}
%================================================================

After nucleation, the daughter undergoes rapid expansion.  Mode~$\ell$ exits the
causal horizon $\ln(\ell/2)$ e-folds after $\ell = 2$, giving:
\begin{equation}
\frac{\text{amplitude at }\ell}{\text{amplitude at }\ell = 2} = \left(\frac{2}{\ell}\right)^{\!2}.
\label{eq:dilution}
\end{equation}

\begin{theorem}[Multipole Dilution, DR]\label{thm:dilution}
The nucleation signal at multipole~$\ell$ is diluted by $(2/\ell)^2$ relative to
$\ell = 2$.  This is parameter-free---it follows from the kinematics of horizon
exit during inflation.
\end{theorem}

Consequences: $\ell = 3$ retains 44\% amplitude (20\% power); $\ell = 4$ retains
25\% amplitude (6\% power); $\ell \geq 6$ is below 1\% power.  For the signal to
be visible, $N_{\mathrm{total}} \approx 64$--$65$ e-folds~[CV].

%================================================================
\section{The CMB Anomalies Explained}\label{sec:cmb}
%================================================================

\subsection{Five anomalies in $\Lambda$CDM}\label{sec:anomalies}

Planck data reveals five anomalies at the largest CMB angular scales:
(1)~$C_2 \approx 200\,\mu\mathrm{K}^2$ vs.\ $\Lambda$CDM $\approx 1100\,\mu\mathrm{K}^2$
(80\% suppression); (2)~$\ell = 2$ and $\ell = 3$ axes aligned to $9$--$13^\circ$;
(3)~N/S hemisphere asymmetry $A_{\mathrm{NS}} \approx 0.07$; (4)~$C_3$ consistent
with $\Lambda$CDM; (5)~$\ell \geq 4$ normal.

\subsection{The RA explanation}\label{sec:explanation}

All five features follow from the Kerr nucleation geometry:
\begin{enumerate}[nosep]
\item The coherent nucleation quadrupole partially cancels the stochastic
  quadrupole, suppressing~$C_2$.
\item Both $\varepsilon_2$ and $\varepsilon_3$ are aligned with the parent's spin
  axis.  The alignment is structural, not coincidental.
\item The $\varepsilon_3 P_3(\cos\theta)$ term is odd under
  $\theta \to \pi - \theta$, creating hemisphere asymmetry.
\item The $(2/\ell)^2$ dilution makes the $\ell = 3$ nucleation signal 44\% of
  $\ell = 2$---too weak to suppress~$C_3$ but sufficient for alignment.
\item The $(2/\ell)^2$ dilution makes the nucleation signal negligible at
  $\ell \geq 4$.
\end{enumerate}

\subsection{Three parameters, zero freedom}\label{sec:threeparam}

The model has three parameters: parent spin~$a_*$, extra e-folds~$\Delta N$, and
accretion asymmetry $\delta_{\mathrm{acc}} = \varepsilon_3/\varepsilon_2 \approx
0.69$.  These are constrained by three observations ($C_2$ suppression, alignment
angle, N/S asymmetry) with two consistency checks ($C_3$ not suppressed, higher
$\ell$ normal).  The system is just-determined: zero remaining freedom.  For any
$a_* \in [0.3, 0.998]$, $\Delta N \approx 3.3$--$5.0$, giving
$N_{\mathrm{total}} \approx 64$--$65$~[CV].

$\Lambda$CDM has no mechanism for any of these anomalies.  RA predicts all five as
structural consequences of Kerr nucleation.

%================================================================
\section{The Cosmic Web}\label{sec:web}
%================================================================

The nucleation perturbations seed structure formation through a feedback mechanism
intrinsic to RA.  The field equations couple local expansion rate inversely to
baryon density: underdense regions expand faster; overdense regions concentrate
further.  This positive feedback drives filamentary cosmic web topology as a
dynamical attractor---inevitable for any universe with density perturbations.

\begin{prediction}[Cosmic web alignment]\label{pred:web}
The large-scale orientation of the cosmic web correlates with the CMB preferred
axis.  Testable with DESI/Euclid.  $\Lambda$CDM predicts no such correlation.
\end{prediction}

%================================================================
\section{The Fate of the Universe}\label{sec:fate}
%================================================================

\subsection{Dual severance mechanism}\label{sec:dual}

The void--filament feedback does not stop.  Both endpoints lead to causal
disconnection:

\emph{Starvation severance.}  Voids approach $\mu = 1$, below which the graph
fragments---quiet disconnection by depletion.

\emph{Density severance.}  Dense boundaries approach $\mu_2 \approx 5.7$, where
$P_{\mathrm{acc}} \to 0$ and new severances fire, nucleating daughter universes.

\subsection{No heat death}\label{sec:noheat}

Heat death requires a single connected component persisting for infinite time.
RA's fragmentation prevents this: every density perturbation drives regions away
from $\mu = 1$, and both directions terminate in disconnection.  The universe does
not equilibrate.  It \emph{fragments}.

%================================================================
\section{The Generational Cycle}\label{sec:cycle}
%================================================================

At dense boundaries, the cycle begins again.  Each daughter inherits initial
conditions from its parent's local structure.  The generational chain:
\begin{quote}
Parent SMBH ($a_*$, accretion) $\to$ anisotropic $\Phi_{\mathrm{boundary}}$ $\to$
daughter initial conditions $\to$ CMB anomalies $\to$ structure formation $\to$
cosmic web $\to$ starvation severance + boundary nucleation $\to$ granddaughter
universes $\to$ cycle continues.
\end{quote}
The physics is identical across generations---only initial conditions vary.  The
``multiverse'' of RA is a multiverse of initial conditions, not of different physics.

%================================================================
\section{Boltzmann Brain Prohibition}\label{sec:bb}
%================================================================

In $\Lambda$CDM with $\Lambda > 0$ and infinite time, Boltzmann Brain nucleations
are inevitable and should vastly outnumber ordinary observers.

RA prohibits them by two mechanisms.  First, \emph{structurally}: the actualization
filter selects based on relative entropy, not Boltzmann statistics; random
fluctuations cannot pass the filter to build ordered complexity.  Second,
\emph{dynamically}: no causal patch persists long enough, because fragmentation
disconnects every component on a finite timescale.

%================================================================
\section{Predictions}\label{sec:predictions}
%================================================================

\begin{prediction}[$\Lambda = 0$]\label{pred:lambda}
Any confirmed nonzero cosmological constant (beyond effective $\Lambda$ from
misidentified density effects) falsifies the framework.
\end{prediction}

\begin{prediction}[No WIMP signal]\label{pred:wimp}
No dark matter signal below the neutrino floor.  LZ and XENONnT results are
consistent.
\end{prediction}

\begin{prediction}[Parent SMBH mass]\label{pred:parent}
$M_{\mathrm{parent}} \approx 2.5$--$6.6 \times 10^6\,M_\odot$.
\end{prediction}

\begin{prediction}[CMB axis persistence]\label{pred:axis}
Quadrupole--octupole alignment persists in CMB-S4 at $> 3\sigma$.
\end{prediction}

\begin{prediction}[N/S asymmetry alignment]\label{pred:ns}
Hemisphere asymmetry aligned with the axis-of-evil axis.
\end{prediction}

\begin{prediction}[Near-minimum inflation]\label{pred:inflation}
$N_{\mathrm{total}} \approx 64$--$65$ e-folds.  Testable via tensor-to-scalar ratio.
\end{prediction}

\begin{prediction}[Cosmic web--CMB correlation]\label{pred:webaxis}
Large-scale cosmic web orientation correlates with CMB preferred axis.
\end{prediction}

\begin{prediction}[Multipole hierarchy]\label{pred:hierarchy}
Higher multipoles ($\ell \geq 4$) show decreasing alignment, following $(2/\ell)^2$.
\end{prediction}

%================================================================
\section{Conclusions}\label{sec:conclusions}
%================================================================

The growing DAG hypothesis resolves five cosmological puzzles: the cosmological
constant, dark matter, the Hubble tension, the CMB anomalies, and the heat death
problem.  Each resolution is structural---a direct consequence of the hypothesis.

The most striking result is the identification of the CMB anomalies as signatures
of the spinning supermassive black hole from which our universe nucleated: the
geometric fingerprint of a Kerr SMBH imprinted on the daughter's initial conditions
and diluted by the parameter-free $(2/\ell)^2$ law.  Three parameters account for
five observations with zero remaining freedom.

The framework predicts a universe that does not die but reproduces---fragmenting
through dual severance, with daughters nucleating at dense boundaries.  The cycle is
eternal, the physics identical across generations, and the unique growth rule ensures
every daughter produces the same Standard Model and the same conditions for
complexity.

Companion papers: coupling constants and particle spectrum (Paper~I~\cite{P1}),
measurement problem and quantum predictions (Paper~II~\cite{P2}),
substrate-independent complexity (Paper~IV~\cite{P4}).  Paper~0~\cite{P0} provides
the narrative overview.

\section*{Acknowledgments}

The formal verification programme uses Lean~4/Mathlib.  AI collaboration (Claude,
Anthropic) contributed to derivation execution and red-team review.

% ─────────────────────────────────────────────────────────────────────────────
\begin{thebibliography}{99}

\bibitem{P1}
J.~F.~Sandeman,
``Paper~I: The Physics of a Self-Consistent Causal Graph,''
Paper~I of this series (2026).

\bibitem{P2}
J.~F.~Sandeman,
``Paper~II: Wave Function Collapse as Irreversible Entropy Production,''
Paper~II of this series (2026).

\bibitem{P4}
J.~F.~Sandeman,
``Paper~IV: The Causal Firewall,''
Paper~IV of this series (2026).

\bibitem{P0}
J.~F.~Sandeman,
``Paper~0: A Biography of the Universe,''
Paper~0 of this series (2026).

\bibitem{BenincasaDowker2010}
D.~M.~T.~Benincasa and F.~Dowker,
``The scalar curvature of a causal set,''
\textit{Phys.\ Rev.\ Lett.}\ \textbf{104}, 181301 (2010).

\bibitem{DJW2010}
B.~Durhuus, T.~Jonsson, and J.~F.~Wheater,
``Random causal trees and the causal set approach to quantum gravity,''
\textit{J.\ Stat.\ Phys.}\ \textbf{139}, 859--881 (2010).

\bibitem{Planck2020}
Planck Collaboration,
``Planck 2018 results.\ VI.\ Cosmological parameters,''
\textit{Astron.\ Astrophys.}\ \textbf{641}, A6 (2020).

\bibitem{DowkerGlaser2013}
F.~Dowker and L.~Glaser,
``Causal set d'Alembertians for various dimensions,''
\textit{Class.\ Quantum Grav.}\ \textbf{30}, 195016 (2013).

\bibitem{Yeats2025}
K.~Yeats,
``Combinatorial interpretation of the coefficients of the causal set
d'Alembertian,''
\textit{Class.\ Quantum Grav.}\ \textbf{42}, 145003 (2025);
arXiv:2412.14036.

\end{thebibliography}

\end{document}
